\section{Proofs and explicit examples}\label{app:proofs}
The first two propositions concern exact representations at fixed geometry. The last two concern exact raw reconstruction functions on their stated domains. None of the proofs assumes that the evaluated forecasting representations satisfy dynamical closure.

\subsection{Score recovery and positivity}
\begin{proof}[Proof of Proposition~\ref{prop:normalizer}]
Write $a_j=(1-T_j^m+\tau)^\omega$. Since $T_j^m\in[0,1]$, $\tau^\omega\leq a_j\leq(1+\tau)^\omega$. All affinities are positive, so their sum $Z_b^m$ and memberships are positive. Taking logarithms in $\phi_j^m=a_j/Z_b^m$ gives
$\log\phi_j^m+\zeta_b^m=\omega\log(1-T_j^m+\tau)$.
Exponentiating and solving for $T_j^m$ proves Equation~\eqref{eq:inverse}. Equation~\eqref{eq:memberships} supplies the converse map. Finally,
\[
 \phi_j^m=\frac{a_j}{a_j+\sum_{i\ne j}a_i}
 \geq\frac{\tau^\omega}{\tau^\omega+(q-1)(1+\tau)^\omega},
\]
because the fraction increases with its numerator affinity and decreases with every competing affinity. The argument applies independently to every partition.
\end{proof}

\subsection{Input recovery and a conditioning bound}
\begin{proof}[Proof of Proposition~\ref{prop:local}]
All memberships are supplied as part of the local encoding. Dividing its $j$th coordinate block by the positive membership gives $P_j^\top(z-\mu_j)$. Adding the known $P_j^\top\mu_j$ and stacking yields $q(z)=\mathcal Pz$. Full column rank makes $\mathcal P^\top\mathcal P$ invertible and gives the displayed recovery formula, hence injectivity.

For the conditioning statement put $w_j=P_j^\top(z-\mu_j)$ and $w'_j=P_j^\top(z'-\mu_j)$. Use $v_j=\phi_jw_j$ and $v'_j=\phi'_jw'_j$, with rule and query arguments suppressed. Then
\[
 w_j-w'_j=\frac{v_j-v'_j}{\phi_j}
             +\frac{\phi'_j-\phi_j}{\phi_j}w'_j.
\]
Stacking these identities and applying the triangle inequality gives
$\|\mathcal P(z-z')\|_2\leq(\|v-v'\|_2+M\|\phi-\phi'\|_2)/\varepsilon$.
The smallest-singular-value inequality
$\|\mathcal P(z-z')\|_2\geq\sigma_{\min}(\mathcal P)\|z-z'\|_2$
proves Equation~\eqref{eq:conditioning}. The constants are assumptions on the query pair or domain; they are not automatically favorable because an inverse exists.
\end{proof}

\subsection{Equality of augmented raw reconstruction spans}
\begin{proof}[Proof of Proposition~\ref{prop:equal}]
A Gaussian Gram matrix is symmetric positive semidefinite. Therefore $\kappa\geq0$ and $C$ is symmetric positive definite. The row-sum assumption gives $C\mathbf1_n=a\mathbf1_n$ with $a=(\kappa+\lambda)^{-1}>0$. The KAHM denominator is $a\mathbf1_n^\top k(z)>0$. Invariance of $\mathcal V$ gives a matrix $B\in\R^{d\times d}$ and a vector $c\in\R^d$ such that $CQ=QB+\mathbf1_nc^\top$. Symmetry then yields
\[
 \mathcal A^K(z)=\frac{(CQ)^\top k(z)}{a\mathbf1_n^\top k(z)}
               =\frac{B^\top\mathcal A^R(z)+c}{a}.
\]
Every coordinate of $g\mathcal A^A=g\mu+P\{gP^\top(z-\mu)\}$ belongs to $\mathcal S_0$. Multiplying the preceding identity by $g$ therefore gives
$\Span(\mathcal S_0,d_K)\subseteq\Span(\mathcal S_0,d_R)$.

The restriction of $C$ to $\mathcal V$ is injective. Because $\mathcal V$ is finite dimensional and $C\mathcal V\subseteq\mathcal V$, it follows that $C\mathcal V=\mathcal V$. Thus there are $D\in\R^{d\times d}$ and $e\in\R^d$ such that $C^{-1}Q=QD+\mathbf1_ne^\top$. Multiplying by $C$ gives $Q=CQD+a\mathbf1_ne^\top$, so
\[
 \mathcal A^R(z)=aD^\top\mathcal A^K(z)+ae.
\]
The same gated-span argument gives the reverse inclusion. No independence of the columns of $Q$, or invertibility of $B$ or $D$, is required. Adjoining the same constant function preserves equality.
\end{proof}

For scalar references $(-1,1)$, let $\rho\in(0,1)$ be the off-diagonal Gaussian Gram entry. With $g=1$, $P=1$, and $\mathcal A^A(z)=z$, direct inversion gives
\[
 \mathcal A^K=c_\rho\mathcal A^R,\qquad
 d_K=c_\rho d_R+(c_\rho-1)z,\qquad
 c_\rho=\frac{1+\rho+\lambda}{1-\rho+\lambda}.
\]
The departure formulas differ, but the augmented spaces containing $z$ coincide. This example concerns jointly variable linear coefficients; freezing a separate common forecast would define a different comparison.

\subsection{Separation on an irregular scalar reference set}
\begin{proof}[Proof of Proposition~\ref{prop:distinct}]
Put $q=(q_1,q_2,q_3)^\top$, $v=C\mathbf1_3$, and $u=Cq$. Symmetry of $C$ gives
$\mathcal A^K=u^\top k/(v^\top k)$ and $\mathcal A^R=q^\top k/(\mathbf1_3^\top k)$.
Since $v_i>0$, the first function is a convex combination of the fixed values $u_i/v_i$, with weights proportional to $v_ik_i(z)$. The second is a convex combination of the references. Both are bounded on $\R$.

Suppose first that $\mathcal A^K=a\mathcal A^R+b$ for constants $a,b$. Let $w=u-bv$. Cross multiplication, with positive denominators, gives
\begin{equation}
 (w^\top k)(\mathbf1_3^\top k)=a(q^\top k)(v^\top k).
 \label{eq:cross}
\end{equation}
For $i\leq j$,
\[
 k_i(z)k_j(z)=e^{-\beta z^2}e^{\beta(q_i+q_j)z}e^{-\beta(q_i^2+q_j^2)/2}.
\]
The six sums $q_i+q_j$ are distinct: strict ordering rules out every possible equality except $q_1+q_3=2q_2$, which is excluded. Exponentials with distinct real rates are linearly independent. To see this, differentiate a proposed zero linear combination at zero through orders zero to five. Its coefficients satisfy a Vandermonde system with nonzero determinant. Multiplication by the common positive factor $e^{-\beta z^2}$ and the nonzero constants preserves independence.

Comparing coefficients in Equation~\eqref{eq:cross} therefore gives $w_i=aq_iv_i$ and $w_i+w_j=a(q_iv_j+q_jv_i)$ for $i<j$. If $a=0$, then $w=0$ and invertibility of $C$ implies $q=b\mathbf1_3$, contradicting distinct references. Otherwise substitution gives $(q_i-q_j)(v_i-v_j)=0$, so all entries of $v$ equal a positive value. The identity $(G+\lambda I_3)v=\mathbf1_3$ would then force constant row sums in $G$.

But its middle row exceeds its first row by $G_{23}-G_{13}>0$ and its last row by $G_{12}-G_{13}>0$: each neighboring reference distance is shorter than $q_3-q_1$. Thus constant row sums are impossible. This rules out an affine relation between the two reconstructions. A relation $\mathcal A^K=a\mathcal A^R+b+cz$ forces $c=0$ by boundedness, so it is also impossible. Finally, $d_m=\mathcal A^m-z$ gives
$\Span(1,z,d_R)=\Span(1,z,\mathcal A^R)$, and membership of $d_K$ in this space would imply membership of $\mathcal A^K$.
\end{proof}

\paragraph{An explicit positive-denominator example.}
Take $q=(-1,0,2)^\top$, $\beta=3/7$, and $\lambda=4$. Write $G=I_3+E$ and $a=5$. The maximum absolute row-sum norm satisfies $\|E\|_\infty=e^{-3/14}+e^{-6/7}<2<a/2$. The convergent Neumann series gives
\[
 \|aC\mathbf1_3-\mathbf1_3\|_\infty
 \leq\frac{\|E\|_\infty/a}{1-\|E\|_\infty/a}<1.
\]
Here the vector infinity norm is the largest absolute coordinate. Consequently all entries of $C\mathbf1_3$ are positive. This is a constructed algebraic example; its chosen regularizer does not assert that the experiment's reference-fitted regularizers satisfy the same condition.
